\documentclass{article}

\usepackage[spanish]{babel}
\usepackage[utf8]{inputenc}
\usepackage[T1]{fontenc}
\usepackage{amsmath,amsfonts,amssymb}
\usepackage{graphicx}
\usepackage{booktabs}
\usepackage{anysize}
\usepackage{setspace}
\usepackage{times}

\marginsize{1.5cm}{1.5cm}{1.5cm}{1.5cm}
\parindent=0mm
\parskip=3mm
\renewcommand{\baselinestretch}{1}

\title{Sistema de visi\'on artificial para conducci\'on aut\'onoma con Jetson Nano}
\author{Juan Jos\'e Prado Neira \\ Javier Rego Bargos}
\date{Mayo 2026}

\begin{document}

\maketitle

\section*{Propuesta de trabajo}

\section{Extracto}

El trabajo consiste en desarrollar un sistema de visi\'on artificial para un coche robot basado en Jetson Nano. El objetivo es que el veh\'iculo pueda seguir un circuito o zona transitable y reaccionar ante cajas u obst\'aculos detectados en la imagen.

La soluci\'on utiliza una red neuronal de segmentaci\'on sem\'antica de tipo U-Net. El modelo recibe una imagen RGB y genera dos m\'ascaras: una para el circuito y otra para los obst\'aculos. A partir de estas m\'ascaras se calcula el error lateral del coche y se generan comandos de control normalizados (\textit{steering} y \textit{throttle}).

El sistema se organiza de forma distribuida: la Jetson captura y env\'ia v\'ideo, mientras que el port\'atil ejecuta la inferencia, calcula la acci\'on de control y devuelve las \'ordenes al veh\'iculo.

\section{Introducci\'on}

La visi\'on artificial es una herramienta fundamental en rob\'otica m\'ovil, especialmente cuando se necesita interpretar el entorno de forma flexible. En este proyecto se emplea segmentaci\'on sem\'antica porque permite obtener informaci\'on p\'ixel a p\'ixel de la escena, lo que resulta adecuado para seguir una zona de circulaci\'on y localizar obst\'aculos.

El uso de una Jetson Nano permite trabajar con una plataforma real de bajo coste, cercana a aplicaciones de rob\'otica educativa y prototipado. El proyecto combina captura de v\'ideo, aprendizaje profundo, comunicaci\'on por red y control de actuadores.

\section{Objetivos}

El objetivo general es construir un prototipo funcional de conducci\'on aut\'onoma basado en visi\'on artificial.

Los objetivos concretos son:

\begin{itemize}
    \item Preparar un dataset propio con im\'agenes del circuito y m\'ascaras de segmentaci\'on.
    \item Entrenar una red neuronal capaz de segmentar circuito y obst\'aculos.
    \item Exportar el modelo a ONNX para su uso en inferencia.
    \item Integrar el flujo de v\'ideo Jetson--port\'atil y el env\'io de comandos port\'atil--Jetson.
    \item Convertir las m\'ascaras de segmentaci\'on en \'ordenes de direcci\'on y velocidad.
    \item Evaluar el sistema mediante m\'etricas de segmentaci\'on y pruebas de conducci\'on.
\end{itemize}

\section{Metodolog\'ia}

Primero se capturan im\'agenes reales del circuito y se anotan con dos clases: \texttt{circuito} y \texttt{caja}. A partir de estas anotaciones se generan m\'ascaras binarias para entrenar el modelo.

Despu\'es se entrena una U-Net peque\~na en PyTorch. La red se dise\~na con dos salidas para resolver en una sola inferencia tanto el seguimiento del circuito como la detecci\'on de obst\'aculos.

Finalmente, el modelo se exporta a ONNX y se integra en el sistema \texttt{socket\_vision}. En ejecuci\'on, la c\'amara de la Jetson env\'ia frames al port\'atil; el port\'atil segmenta la imagen, calcula el centro del circuito y env\'ia comandos de control al coche.

\section{Apartado t\'ecnico}

\subsection{Dataset}

\begin{table}[h]
\centering
\begin{tabular}{ll}
\toprule
\textbf{Elemento} & \textbf{Valor} \\
\midrule
Im\'agenes anotadas & 57 \\
Anotaciones totales & 100 \\
Clases & \texttt{circuito}, \texttt{caja} \\
Resoluci\'on original & $640 \times 360$ \\
Im\'agenes de entrenamiento & 45 \\
Im\'agenes de validaci\'on & 12 \\
Formato de m\'ascara & Binaria por clase \\
\bottomrule
\end{tabular}
\end{table}

\subsection{Modelo neuronal}

\begin{table}[h]
\centering
\begin{tabular}{ll}
\toprule
\textbf{Caracter\'istica} & \textbf{Valor} \\
\midrule
Arquitectura & U-Net ligera \\
Entrada & $3 \times 256 \times 256$ \\
Salida & $2 \times 256 \times 256$ \\
Canales de salida & Circuito y obst\'aculo \\
Canales del encoder & 32, 64, 128, 256 \\
Canales del cuello de botella & 512 \\
Etapas de bajada/subida & 4 / 4 \\
Convoluciones 2D & 19 \\
Convoluciones transpuestas & 4 \\
BatchNorm + ReLU & 18 + 18 \\
Par\'ametros entrenables & 7.763.074 \\
\bottomrule
\end{tabular}
\end{table}

\subsection{Entrenamiento e inferencia}

\begin{table}[h]
\centering
\begin{tabular}{ll}
\toprule
\textbf{Elemento} & \textbf{Valor} \\
\midrule
Framework de entrenamiento & PyTorch 2.11 \\
Optimizador & Adam \\
\textit{Learning rate} & $10^{-3}$ \\
\textit{Batch size} & 8 \\
Funci\'on de p\'erdida & BCEWithLogits + Dice Loss \\
M\'etricas & Dice e IoU \\
Checkpoint seleccionado & \texttt{best\_model.pth} \\
Mejor \'epoca & 78 \\
Dice de validaci\'on del mejor modelo & 0.9519 \\
P\'erdida de validaci\'on del mejor modelo & 0.1175 \\
Formato de exportaci\'on & ONNX \\
Resoluci\'on de captura & $512 \times 288$ \\
Frecuencia objetivo & 12 FPS \\
Comandos de salida & \textit{steering}, \textit{throttle} \\
\bottomrule
\end{tabular}
\end{table}

\section{Plan de trabajo}

\begin{itemize}
    \item \textbf{Fase 1:} revisi\'on del problema, definici\'on de requisitos y preparaci\'on del entorno.
    \item \textbf{Fase 2:} captura, anotaci\'on y organizaci\'on del dataset.
    \item \textbf{Fase 3:} entrenamiento, validaci\'on y exportaci\'on del modelo.
    \item \textbf{Fase 4:} integraci\'on del modelo en el flujo Jetson--port\'atil.
    \item \textbf{Fase 5:} ajuste del control aut\'onomo y pruebas con obst\'aculos.
    \item \textbf{Fase 6:} evaluaci\'on final, an\'alisis de resultados y redacci\'on de la memoria.
\end{itemize}

\section{Evaluaci\'on prevista}

La evaluaci\'on del sistema completo incluir\'a FPS, latencia, estabilidad del seguimiento del circuito, respuesta ante obst\'aculos y necesidad de intervenci\'on manual.

\end{document}
